\section{结论与后续工作}
\label{sec:conclusion}

在一个固定的前沿模型上，一个做得好的框架能追上服务商自家的智能体，并让它更省，但还没有超过它：我们在
Terminal-Bench~4.1 上与 GPT-6.2 Sol、max 强度的 Codex 持平，输入 token 少约 30\%。

本文有三个主要结论。上下文工程的回报在成本，不在准确率：只追加的提示词把成本降低三分之一到一半，之后的手段又降低
7\% 到 20\%。结构只在工作可以拆分的地方有回报：在单个上下文装得下的任务上，一个能力强的负责者胜过所有团队，同一模型
的第二意见只会确认共同的误读，抓不住它。判断跟随模型：提示、规则和人设改变了模型下一步做什么，却没有改变它选择什么；
摆到它面前的事实和测量结果才改变了选择。

因此，要拉开差距，需要和我们已经放弃的那些手段不同类型的手段：
\begin{enumerate}[leftmargin=1.6em]
  \item \textbf{证据，而不是建议。}在决策点上，把数据的测量结果和每种待定解读的可执行后果摆到领队面前；这是唯一
    改变过选择的通道。
  \item \textbf{不共享同一解读的多样性。}同一模型的投票者在错误的解读上达成了一致。维护者榜单的第一名是一个多模型
    智能体（71.2\%），而在 Codex 中 GPT-6.2 Astra 比 GPT-6.2 Sol 高 5.5 个百分点；在解读或选择这一步加入第二个 GPT 模型，
    是显而易见的下一个实验。
  \item \textbf{更多次运行，加上有依据的选择者。}在 GPT-6.2 上，一个完美的选择者从两次运行中挑选，可以在 30 道题上多出
    约 1.6 道。选择者无论用 GPT 还是 Claude，在明显缺陷上都选得对，在需要判断的点上都选错（28 次中 5 次），所以选择者
    必须依据执行过的检查，而不是意见。
  \item \textbf{可比较的报告。}用一台有 GPU 的机器跑全部 66 道题、每题五次尝试，就能与维护者的协议一致；这是向第三方
    排行榜提交的门槛。
\end{enumerate}

\section*{数据与代码可用性}
Harbor 适配器以及复现最终运行的完整命令已经公开 \citep{zhang2026adapter}（附录~\ref{app:repro}）；已发布的服务端按
\citep{zeroth2026server} 中的文档安装。Terminal-Bench~4.1 运行的每次尝试结果（奖励、智能体用时、token 和按公开标价
计算的成本）以及可读题的轨迹由作者保存。30 道留出题的轨迹仍未读取，以便这个集合今后还能用于确认。

\section*{AI 辅助说明}
本文的实验在一个 AI 编程助手（Claude Code，模型为 Claude Opus 5.5）的帮助下设计、运行和分析，论文初稿也由它在作者
指导下撰写。作者对内容负责。

\section*{利益冲突}
作者创办了 The Zeroth，该公司开发 The One。
